\chapter{2018:Peter Scholze}\label{ch2018peterscholze:peter-scholze-fields-medal-2018-a-comprehensive-report-on-his-work-and-impact}

\begin{center}\rule{0.5\linewidth}{0.5pt}\end{center}

\textbf{Prepared:} July 2026

\textbf{Subject:} Peter Scholze (b.\ December 11, 1987, Dresden, Germany)

\textbf{Position:} Director, Max Planck Institute for Mathematics, Bonn; Professor, University of Bonn

\textbf{Doctoral Advisor:} Michael Rapoport (University of Bonn)

\textbf{Fields Medal Citation:} \emph{``For transforming arithmetic algebraic geometry over p-adic fields through his introduction of perfectoid spaces, with application to Galois representations, and for the development of new cohomology theories.''}

\begin{center}\rule{0.5\linewidth}{0.5pt}\end{center}

\begin{prereqs}
Algebraic geometry (schemes, cohomology), number theory (p-adic numbers, local fields, Galois representations), basic representation theory, commutative algebra.
\end{prereqs}

\begin{objectives}
Explain perfectoid spaces and the tilting correspondence; describe p-adic Hodge theory and its extension to rigid-analytic varieties; understand the p-adic Langlands correspondence and prismatic cohomology.
\end{objectives}

\section{Problem History}

p-adic geometry studies geometric objects defined over p-adic fields $\mathbb{Q}_p$. In the 1960s, Tate introduced rigid-analytic geometry as a p-adic analogue of complex analytic geometry. p-adic Hodge theory, developed by Fontaine, Faltings, and others in the 1970s--1990s, establishes comparison theorems between different p-adic cohomology theories: \'{e}tale, de Rham, and crystalline cohomology. These comparison theorems are fundamental tools in arithmetic geometry, connecting the algebraic and analytic aspects of p-adic varieties.

Faltings' almost purity theorem was a key technical result, but its proof was extremely complex and difficult to generalize. The \textbf{weight-monodromy conjecture}, formulated by Deligne in the 1970s, predicts that the weight filtration on \'{e}tale cohomology coincides with the monodromy filtration for proper smooth varieties over local fields. This conjecture, one of the deepest open problems in arithmetic geometry, encodes deep information about the geometry of varieties over local fields.

The central difficulty was the gap between \textbf{mixed characteristic} (characteristic 0 with residue field of characteristic $p$) and \textbf{equal characteristic} $p$. Many results are known in equal characteristic but remain conjectural in mixed characteristic. Bridging this gap had been an elusive goal in arithmetic geometry for decades. The classical tools of p-adic geometry---rigid-analytic spaces, formal schemes, and $p$-divisible groups---were not sufficient to overcome this barrier.

The \textbf{p-adic Langlands correspondence}, formulated by Langlands and Breuil, predicts a deep connection between Galois representations of $\mathbb{Q}_p$ and automorphic representations of $\operatorname{GL}_n(\mathbb{Q}_p)$. The construction of the $p$-adic local Langlands correspondence requires a detailed understanding of the cohomology of $p$-adic local systems on p-adic analytic spaces, which was not available before Scholze's work.

The \textbf{Fontaine--Wintberger conjecture} (1985) predicts that the semisimple part of the monodromy operator on the \'{e}tale cohomology of a proper smooth variety over a p-adic field can be computed from the Galois action on the crystalline cohomology. This conjecture, closely related to the weight-monodromy conjecture, was known only in special cases before the development of perfectoid spaces.

The \textbf{Fargues--Fontaine curve} is a geometric object that parametrizes $p$-divisible groups and provides a geometric framework for p-adic Hodge theory. The curve is a Dedekind domain whose points correspond to perfectoid fields, and its cohomology theory encodes the comparison theorems of p-adic Hodge theory. Understanding the geometry of the Fargues--Fontaine curve was a major open problem.

\section{How Was It Solved}

\textbf{Perfectoid fields and tilting (2012).} A \textbf{perfectoid field} is a complete topological field $K$ with non-discrete valuation where the Frobenius map $\Phi: \mathcal{O}_K/(p) \to \mathcal{O}_K/(p)$ is surjective. The \textbf{tilting} operation associates to a perfectoid field $K$ of characteristic 0 a perfectoid field $K^\flat$ of characteristic $p$ via the inverse limit $K^\flat = \varprojlim_{x \mapsto x^p} K$. For example, the tilt of $\mathbb{Q}_p(p^{1/p^\infty})$ is $\mathbb{F}_p((t))$. This provides a ``bridge'' between the two characteristics, reducing problems in mixed characteristic to equal characteristic where tools are more developed.

\textbf{Perfectoid spaces (2012--2013).} A perfectoid space is an adic space covered by affinoid perfectoid algebras. The tilting correspondence extends to geometry: there is an equivalence of categories between perfectoid spaces over $K$ and over $K^\flat$. This equivalence preserves many geometric properties, including cohomology groups. The key theorem is that the \'{e}tale cohomology of a perfectoid space over $K$ is isomorphic to the \'{e}tale cohomology of its tilt over $K^\flat$.

The theory of perfectoid algebras provides the algebraic foundation. A perfectoid algebra over $K$ is a complete $p$-adically completed tensor product $R = W(\mathcal{O}_{K^\flat})[\![x_1^{1/p^\infty}, \ldots, x_n^{1/p^\infty}]\!]/I$. The tilt of $R$ is the perfectoid $\mathbb{F}_p$-algebra $R^\flat = \varprojlim_{x \mapsto x^p} R/pR$. The equivalence of categories between perfectoid $K$-algebras and perfectoid $K^\flat$-algebras is the algebraic content of the tilting correspondence.

\textbf{Weight-monodromy conjecture (2012).} Using perfectoid spaces, Scholze proved the weight-monodromy conjecture for semistable varieties. The proof reduces the problem to the equal characteristic case using the tilting correspondence, where the conjecture follows from the work of Gabber and de Jong. This was the first major application of perfectoid spaces and demonstrated their power.

The weight-monodromy conjecture predicts that for a proper smooth variety $X$ over a local field $K$ with residue field of characteristic $p > 0$, the weight filtration on \'{e}tale cohomology $H^i_{\text{\'{e}t}}(X_{\overline{K}}, \mathbb{Q}_p)$ coincides with the monodromy filtration. The weight filtration is defined using the eigenvalues of the Frobenius automorphism, and the monodromy filtration is defined using the nilpotent operator $N$ coming from the monodromy of the Galois representation. The conjecture predicts that these two filtrations are the same, which is a deep arithmetic property of the variety.

The proof using perfectoid spaces proceeds by reducing to the equal characteristic case, where the conjecture is known by the work of Gabber and de Jong. The key step is to lift the variety from mixed characteristic to equal characteristic using the tilting correspondence, apply the known result in equal characteristic, and then transfer back to mixed characteristic. This argument demonstrates the power of perfectoid spaces as a tool for transferring results between different characteristics.

\textbf{p-adic Hodge theory for rigid-analytic varieties (2013).} Scholze extended p-adic Hodge theory to rigid-analytic varieties, proving the Hodge--Tate decomposition. The proof uses perfectoid spaces to reduce the problem to the study of sheaves on the perfectoid tower, which can be analyzed using the theory of $\varprojlim$-sheaves.

\textbf{Prismatic cohomology (2018--2020).} With Bhatt, Scholze developed prismatic cohomology, a new cohomology theory for rings equipped with a prime element $p$. Prismatic cohomology provides a single framework that specializes to \'{e}tale, de Rham, and crystalline cohomology, unifying the main p-adic cohomology theories. The theory is defined using the language of \textbf{prisms}, which are certain complete cotangent complexes equipped with a compatible Frobenius lift.

The key innovation of prismatic cohomology is that it provides a functorial construction that interpolates between the three classical p-adic cohomology theories. The prismatic site of a ring $R$ is a category of prisms over $R$ equipped with a Grothendieck topology, and the prismatic cohomology is the cohomology of the structure sheaf on this site. The theory has a natural filtration whose graded pieces recover the classical theories.

A \textbf{prism} is a pair $(A, I)$ where $A$ is a $\delta$-ring (a ring equipped with a derivation $\delta: A \to A$ satisfying certain axioms) and $I \subset A$ is an ideal such that $A/I$ is $\delta$-closed and the $p$-adic completion $\widehat{A}$ is faithfully flat over $\mathbb{Z}_p$. The prism provides a framework for lifting the Frobenius morphism from characteristic $p$ to characteristic 0, which is the key operation in p-adic Hodge theory. The prismatic cohomology of a ring $R$ is defined as the cohomology of the structure sheaf on the prismatic site, which is a category of prisms over $R$ equipped with a Grothendieck topology.

The relationship between prismatic cohomology and the classical theories is given by specialization functors:
\[
H^*_{\mathrm{prism}}(R) \otimes_A \widehat{A}[1/p] \cong H^*_{\mathrm{dR}}(R)
\]
for the de Rham comparison, and similar comparisons for crystalline and \'{e}tale cohomology. These comparisons show that prismatic cohomology is a true generalization of the classical theories, providing a single framework that unifies them.

\section{Impact of This Discovery}

Perfectoid spaces have transformed arithmetic geometry. They provide a clean framework for Faltings' almost purity theorem, enable the study of infinite-level Shimura varieties (which turn out to be perfectoid spaces), and power the Fargues--Scholze theorem geometrizing the local Langlands correspondence via the Fargues--Fontaine curve.

Prismatic cohomology, described as a major step toward motivic cohomology, provides a single theory specializing to all three classical p-adic cohomology theories. This unification has led to new results on the structure of p-adic cohomology and its applications to number theory.

With Clausen, Scholze launched condensed mathematics, resolving the non-abelian nature of topological groups and providing new foundations for functional analysis. The formalization of perfectoid spaces in the Lean proof assistant confirmed that these complex definitions are fully rigorous and computationally tractable.

The \textbf{Fargues--Scholze theorem} geometrizes the local Langlands correspondence by constructing a sheaf on the Fargues--Fontaine curve that encodes the relationship between Galois representations and automorphic representations. This theorem has been influential in the study of the $p$-adic Langlands program and has connections to the theory of diamond morphisms.

The \textbf{Fargues--Fontaine curve} is a geometric object that parametrizes $p$-divisible groups and provides a geometric framework for p-adic Hodge theory. For a perfectoid field $K$, the Fargues--Fontaine curve $\mathbb{B}_{\mathrm{FF}}$ is a Dedekind domain whose points correspond to perfectoid fields over $K$. The cohomology of vector bundles on $\mathbb{B}_{\mathrm{FF}}$ encodes the comparison theorems of p-adic Hodge theory: the de Rham cohomology of a variety over $K$ is related to the \'{e}tale cohomology of its cohomology sheaf on $\mathbb{B}_{\mathrm{FF}}$.

The Fargues--Fontaine curve has a rich geometry that connects p-adic Hodge theory to the theory of vector bundles on curves. The classification of vector bundles on $\mathbb{B}_{\mathrm{FF}}$ is analogous to the Harder--Narasimhan filtration for vector bundles on algebraic curves, and the slope of a vector bundle determines its splitting behavior. This connection between p-adic Hodge theory and the geometry of curves is one of the most surprising aspects of Scholze's work.

\section{Modern Views and Improvements}

Scholze and Clausen are extending condensed mathematics to complex geometry, providing a unified treatment of real, p-adic, and complex analytic geometry. Their six-functor formalism for analytic geometry aims to establish Grothendieck's framework in the perfectoid setting. The theory of condensed sets, which replaces the category of topological spaces, resolves several foundational issues in functional analysis and algebraic topology.

The geometrization of the local Langlands correspondence, developed with Fargues, continues to be refined. Proving full faithfulness of the geometric construction and extending to general reductive groups are major open challenges that remain active areas of research. The connections to the theory of diamonds and to the Langlands program for $\operatorname{GL}_n$ over local fields are being actively developed.

Prismatic cohomology remains an active area: extending the theory to more general bases, proving Poincar\'e duality and K\"unneth formulas, and applications to special values of L-functions are current frontiers. The connections to condensed mathematics and to the theory of derived algebraic geometry are being explored.

The theory of prismatic cohomology has led to new results on the structure of p-adic cohomology. For example, the prismatic cohomology of a smooth scheme over $\mathbb{Z}_p$ has a natural filtration whose graded pieces are the de Rham cohomology, the crystalline cohomology, and the \'{e}tale cohomology. This filtration provides a unified framework for understanding the relationships between these classical theories and has led to new results on the structure of p-adic cohomology.

The \textbf{tilting equivalence} has been extended to other geometric settings, including the study of Shimura varieties, the construction of $p$-adic period domains, and the theory of adic spaces over non-archimedean fields. The tilting equivalence provides a systematic method for transferring results from equal characteristic to mixed characteristic, and vice versa.

The development of condensed mathematics, in collaboration with Dustin Clausen, represents a foundational shift in the approach to functional analysis and algebraic topology. By replacing topological spaces with condensed sets, Scholze and Clausen resolve several foundational issues, including the non-abelian nature of topological groups and the failure of the tensor product to preserve topological properties. The condensed framework provides new foundations for p-adic Hodge theory and has potential applications to the study of perfectoid spaces and prismatic cohomology.

The prismatic cohomology developed by Scholze and Bhatt has been extended to more general settings. The theory has been applied to the study of formal groups, to the theory of p-divisible groups, and to the construction of new cohomology theories for algebraic varieties. The connections between prismatic cohomology and other cohomology theories, such as motivic cohomology and syntomic cohomology, are being actively explored. These connections have the potential to unify the main p-adic cohomology theories and to provide new tools for studying the arithmetic of algebraic varieties.

The theory of diamonds, developed by Scholze as a framework for the geometry of the Fargues--Fontaine curve, has become a central tool in p-adic geometry. Diamonds are a class of analytic stacks that provide a geometric framework for the study of perfectoid spaces and their moduli. The theory of diamonds has been applied to the study of Shimura varieties, the construction of p-adic period domains, and the geometrization of the local Langlands correspondence. The development of the theory of diamonds continues to be an active area of research.

The Fargues-Scholze theorem on the local Langlands correspondence has been extended to other reductive groups. The geometrization of the local Langlands correspondence constructs a sheaf on the Fargues-Fontaine curve that encodes the relationship between Galois representations and automorphic representations. The construction uses perfectoid spaces and the theory of diamonds, and the resulting sheaf provides a geometric framework for the local Langlands program. The extension to general reductive groups is an active area of research that builds on Scholze's foundational work.

The study of Shimura varieties using perfectoid spaces has led to new results on the arithmetic of these fundamental objects. The infinite-level Shimura varieties, which parametrize abelian varieties with additional structure at all levels, turn out to be perfectoid spaces. This observation, made by Scholze, has implications for the study of Galois representations associated to automorphic forms and for the construction of p-adic L-functions. The connection between perfectoid spaces and Shimura varieties is an active area of research that builds on Scholze's foundational work.

The theory of condensed mathematics has potential applications to the foundations of algebraic geometry. The traditional foundations of algebraic geometry, based on schemes and topological spaces, have limitations when dealing with analytic objects such as complex manifolds and p-adic analytic spaces. The condensed framework provides a unified treatment of algebraic and analytic geometry, which has potential applications to the study of mixed-characteristic geometry and to the development of new foundations for algebraic geometry.

The study of p-adic integration and p-adic Hodge theory using prismatic cohomology is an active area. The p-adic integral, which generalizes the classical integral to the p-adic setting, has applications to the study of p-adic L-functions and to the theory of p-adic modular forms. The development of new techniques for computing p-adic integrals using prismatic cohomology, building on Scholze's foundational work, is an active area that connects p-adic Hodge theory to analytic number theory.

\begin{center}\rule{0.5\linewidth}{0.5pt}\end{center}

\section*{Summary}
\addcontentsline{toc}{section}{Summary}

Scholze created perfectoid spaces, providing a bridge between characteristic 0 and characteristic $p$. His work on p-adic Hodge theory, prismatic cohomology, and condensed mathematics has reshaped arithmetic geometry and established new foundations for studying $p$-adic fields.
